\documentclass[10pt]{article}
 
% basic packages
\usepackage[margin=1in]{geometry}
\usepackage[pdftex]{graphicx}
\usepackage{amsmath,amssymb,amsthm}
\usepackage{custom}
\usepackage{comment}
\usepackage{enumitem}

% page formatting
\usepackage{fancyhdr}
\pagestyle{fancy}

\begin{document}
\title{Reading List for 2026}
\date{Last Updated On: Feb 27th 2026}
\maketitle

\fancyhf{}
\fancyhead[L]{2026 Reading List}
\fancyhead[R]{\leftmark}
\fancyhead[C]{\thepage}

% Making sure the right header is NOT in all caps %
\renewcommand{\sectionmark}[1]{\markright{\textsf{\arabic{section}: #1}}}
\rhead{Section\ \nouppercase{\rightmark}}
\chead{}
\cfoot{}

\begin{comment}
   ------ THESE ARE THE PREVIOUS COMMANDS FOR THE HEADER, FOOTER, ETC -------
\renewcommand{\sectionmark}[1]{\markright{\textsf{\arabic{section}. #1}}}
\renewcommand{\subsectionmark}[1]{}
\lhead{\scriptsize\textbf{\thepage}\ \ \nouppercase{\rightmark}}
\chead{}
\rhead{\normalfont MAP2302 - Spring 2026}
\lfoot{}
\cfoot{}
\rfoot{}
\setlength{\headheight}{5pt}
   ------ THESE ARE THE PREVIOUS COMMANDS FOR THE HEADER, FOOTER, ETC -------
\end{comment}

\linespread{1.03} % give a little extra room
\setlength{\parindent}{0pt}
\setlength{\parskip}{6pt}
\setcounter{secnumdepth}{2} % Numbered subsections  
\setcounter{tocdepth}{2} % numbered subsubsections in ToC

\tableofcontents

%%%%% COMPUTER SCIENCE LITERATURE %%%%%
\section{Computer Science}
This section encompasses my reading in Computer Science, 
spanning both technical and narrative works. Topics include, but are not limited to:
\begin{enumerate}
    \item Collected Works 
    \item Fiction and biographical accounts of notable and up and coming computer scientists
    \item Conceptual textbooks ranging from foundational to advanced topics, such as: 
    \begin{enumerate}
        \item \textbf{Data Structures \& Algorithms }
                \begin{enumerate}
                    \item \underline{\textbf{Data Structures}} can be defined as the medium through which data and is stored to be 
                        accessed later on. Simply put, a data structure is a collection of data values, the relationships that take place among them, and 
                        the functions or operations that can be applied to the data
                \end{enumerate}
        \item \textbf{Computer Organization and Design } \\
        \item \textbf{Operating Systems} \\
        \item \textbf{Discrete Structures} \\
        \item \textbf{Machine Learning \& Neural Networks } \\
        \item \textbf{Reinforcement Learning } \\
        \item \textbf{Natural Language processing (NLP)} \\
        \item \textbf{Computer Vision } \\
        \item \textbf{Information Theory} \\
    \end{enumerate}
\end{enumerate}


\subsection{Fiction}
\subsection{Nonfiction}
\subsection{Research Literature}
\subsection{Textbooks}


%%%%% LEISURE LITERATURE %%%%%
\section{Leisure}
\subsection{Fiction}
\subsection{Nonfiction}
\subsection{Research Literature}
\subsection{Textbooks}


%%%%% MATHEMATICS LITERATURE %%%%%
\section{Mathematics}
\subsection{Fiction}
\subsection{Nonfiction}
\subsection{Research Literature}
\subsection{Textbooks}


%%%%% PHYSICS LITERATURE %%%%%
\section{Physics}
\subsection{Fiction}
\subsection{Nonfiction}
\subsection{Research Literature}
\subsection{Textbooks}


\section{Books Completed w/ Time Of Completion}
\subsection{Computer Science}
\subsection{Leisure}
\subsection{Mathematics}
\subsection{Physics}
\newpage

\end{document}